\section{系统性总结}
\presenter{总结讲解人待填写}
\source{本部分为结构模板；待六篇个人讲解完成后共同填写}
\begin{frame}[plain]
    \vspace{9mm}{\small\color{accent}共同讨论}\par
    \vspace{9mm}{\fontsize{28}{34}\selectfont\bfseries 系统性总结}\par
    \vspace{5mm}{\Large 从六篇工作形成统一认识}\par
    \vspace{7mm}\textcolor{accent}{\rule{\textwidth}{0.8pt}}\par
    \vspace{3mm}讲解人待填写\hfill 10--15 分钟
    \note{以下保留总结页面结构。待大家完成个人部分并讨论后填写，不提前写成已形成的小组共识。默认本部分约 12 分钟。}
\end{frame}

\begin{frame}{统一比较：研究对象、机制与边界}
    \small
    \begin{tabularx}{\linewidth}{@{}p{27mm}YYY@{}}
        \toprule
        工作 & 记忆表示与对象 & 核心干预环节 & 适用条件与局限 \\
        \midrule
        StructMem & 待共同整理 & 待共同整理 & 待共同整理 \\
        Hindsight & 待共同整理 & 待共同整理 & 待共同整理 \\
        Semantic XPath & 待共同整理 & 待共同整理 & 待共同整理 \\
        HyMEM & 待共同整理 & 待共同整理 & 待共同整理 \\
        MetaMem & 待共同整理 & 待共同整理 & 待共同整理 \\
        Text2Mem & 待共同整理 & 待共同整理 & 待共同整理 \\
        \bottomrule
    \end{tabularx}
    \aside{填写各论文实际处理的问题，避免把不同研究层次视为可直接替换的方法。}
    \note{建议约 2 分钟。先说明统一坐标，再解释主要差异；不要逐格朗读。}
\end{frame}

\begin{frame}{机制之间的联系与取舍}
    \begin{columns}[T,onlytextwidth]
        \column{.47\textwidth}
        \textbf{可能互补的环节}\par\vspace{4mm}
        \begin{itemize}\setlength{\itemsep}{3mm}
            \item 待填写：哪些问题需要联合处理？
            \item 待填写：两种机制的接口是什么？
            \item 待填写：原文或实验依据是什么？
        \end{itemize}
        \column{.47\textwidth}
        \textbf{组合的条件与代价}\par\vspace{4mm}
        \begin{itemize}\setlength{\itemsep}{3mm}
            \item 待填写：表示与信息边界是否一致？
            \item 待填写：增加了什么成本或风险？
            \item 待填写：哪些效果还需要验证？
        \end{itemize}
    \end{columns}
    \note{建议约 2 分钟。选择一到两组有依据的联系讲清楚，不要求把六篇强行接成一个系统。}
\end{frame}

\begin{frame}{共同案例：不同方法分别解释什么}
    \textbf{共同场景与问题}\par\vspace{3mm}
    待讨论后选定：给出同一段历史、同一个查询，以及需要保留的关键证据。\par
    \vspace{5mm}
    \begin{tabularx}{\linewidth}{@{}p{28mm}Y@{}}
        \toprule
        比较内容 & 待填写的分析 \\
        \midrule
        组织与检索 & 记忆如何形成，哪些证据能被取回？ \\
        更新与使用 & 新旧信息如何处理，回答或操作如何形成？ \\
        解释边界 & 哪些是论文直接支持，哪些是推演或尚待验证？ \\
        \bottomrule
    \end{tabularx}
    \note{建议约 2 分钟。明确案例来自论文、自拟讨论还是实际实验，不把推演结果当作已经运行的系统对照。}
\end{frame}

\begin{frame}{实验结果的可比条件}
    \begin{tabularx}{\linewidth}{@{}p{32mm}Y@{}}
        \toprule
        条件 & 共同核对后填写 \\
        \midrule
        任务与数据 & 数据划分、题型与有效样本是否一致？ \\
        模型与信息 & 基础模型、优化反馈、测试可用信息有何差异？ \\
        检索与预算 & 上下文数量、token 预算与构建成本如何计量？ \\
        指标与评价 & 各指标测什么，裁判与评价提示是否一致？ \\
        \bottomrule
    \end{tabularx}\par
    \vspace{4mm}
    \key{待填写：哪些结果可以比较，哪些只能分别解释}
    \note{建议约 2 分钟。重点说明影响判断的条件，避免把不同任务的原论文总分直接排名。}
\end{frame}

\begin{frame}{小组结论与课设选择}
    \begin{enumerate}\setlength{\itemsep}{5mm}
        \item \textbf{共同认识：}待填写一项有依据的结论。
        \item \textbf{研究取舍：}待填写主线选择，以及借鉴或暂不采用的机制。
        \item \textbf{仍待验证：}待填写关键分歧、失败情形和后续验证方式。
    \end{enumerate}
    \note{建议约 2 分钟。结论以实际团队讨论和证据为依据；说明选择理由，不重复六位成员的结尾。}
\end{frame}

\source{}
\begin{frame}[plain]
    \vspace{15mm}{\fontsize{28}{34}\selectfont\bfseries 讨论与问题}\par
    \vspace{7mm}{\Large 长上下文场景结构化记忆建模}\par
    \vspace{6mm}{\color{muted}第 8 队}
    \note{问答时间是否计入总时长，按实际安排处理。}
\end{frame}
